% Generated draft: requires mathematical and source review.
\documentclass[11pt]{article}
\usepackage[margin=25mm]{geometry}
\usepackage{fontspec}
\setmainfont{Latin Modern Roman}
\usepackage{amsmath,amssymb,mathtools,unicode-math}
\setmathfont{Latin Modern Math}
\usepackage{enumitem,needspace,xurl,hyperref}
\setlength{\parindent}{0pt}
\setlength{\parskip}{5pt}
\setlength{\emergencystretch}{3em}
\newcommand{\N}{\mathbb{N}}
\newcommand{\Z}{\mathbb{Z}}
\newcommand{\Q}{\mathbb{Q}}
\newcommand{\R}{\mathbb{R}}
\newcommand{\C}{\mathbb{C}}
\newcommand{\F}{\mathbb{F}}
\newcommand{\A}{\mathbb{A}}
\newcommand{\PP}{\mathbb{P}}
\newcommand{\E}{\mathbb{E}}
\newcommand{\eps}{\varepsilon}
\newcommand{\dd}{\,\mathrm{d}}
\begin{document}
\section*{UnsolvedMath ID: 2691}
\textbf{Source problem code:} catalogue entry
\subsection{Definitions and mathematical statement}
Let $L\subset\R^3$ be a smooth link, and let $y$ denote the positive coordinate in $\R_+=(0,\infty)$. The source proposes a gauge-theoretic construction with boundary at $y=0$ and a special link singularity along $L\times\{0\}$. It poses three connected but distinct tasks.
\begin{itemize}
\item (a) On the four-dimensional space $\R^3\times\R_+$, construct and analyze a moduli theory for solutions of the Kapustin--Witten (KW) equations with Nahm pole boundary conditions, modified along the link. Define appropriate graded solution counts that recover the coefficients of the Jones polynomial of $L$.
\item (b) On $\R_t\times\R^3\times\R_+$, construct the corresponding Haydys--Witten (HW) theory. The source's proposed complex has KW solutions as generators and a differential given by counting HW solutions with appropriate limits as $t\to\pm\infty$ and the link/Nahm pole boundary conditions. The task includes defining the counts and gradings, proving that the differential squares to zero, proving independence from auxiliary choices, and identifying its homology with the chosen version of Khovanov homology of $L$.
\item (c) Replace $\R^3$ by another $3$-manifold $Y$ and use KW or HW equations to define link invariants for $L\subset Y$. The source suggests KW counts might recover classes in the Kauffman bracket skein module or relate to physicists' $\widetilde Z$ invariants, while HW counts might give Khovanov-type homology. Comparisons with Khovanov lasagna skein modules are part of the proposed programme, not an established identification.
\end{itemize}
Khovanov homology is a combinatorially defined link homology whose graded Euler characteristic recovers the Jones polynomial with corresponding conventions. Its cobordism behavior motivates the gauge-theoretic interpretation. In the source, KW and HW refer to the equation systems introduced in [KW07], [Wit12], and [Hay15], and the relevant singular boundary conditions are referred to [Wit12], [MW14], and [MW20]. The complete PDE systems, structure group and bundle, perturbations, behavior at spatial and $y$ infinity, compactification, orientations, gradings, coefficient ring, and reduced/unreduced and Jones normalization choices are not given in this record. These choices are necessary for a single rigorous counting theorem; none is silently selected. Formal solution counts are requirements to construct, not assumed finite or already invariant.

The literature notes credit Witten and Gaiotto--Witten with proposing/sketching the Jones and Khovanov interpretations. They list analytic progress in [MW14], [MW20], [Tau13], [LT20], [He19], [HM19a], [Tau18], [Tau19]. They emphasize that some compactness results, including the cited [Tau18, Theorem B], require positive Ricci curvature of the $3$-manifold, suggesting this as an initial setting for (c), not a theorem for every $Y$. The notes list concrete combinatorial theories for links in $S^3$, interval bundles over surfaces, $\R\mathrm P^3$, and connected sums of $S^1\times S^2$, with references [APS04], [Man07], [Gab13], [Roz10], [Wil21]. Their scope is reported as source context rather than an independently updated exhaustive classification. The $\widetilde Z$ direction refers to source Problem 3.68 and [GPV17], [GPPV20]; [MWW22] supplies the lasagna comparison context.

The supplied triage checked 2026-08-17 labels the three-part programme open (OPEN-TRIAGE) and says no complete analytic construction matching all requested properties was verified. The registry is also open. The intended task is rigorous gauge-theoretic recovery, not merely naming a proposed physics interpretation.
\subsection{Short English statement}
Make the proposed Kapustin--Witten counts recover the Jones polynomial, and the Haydys--Witten complex recover Khovanov homology, with link-sensitive Nahm pole boundaries. Then extend the construction to links in other $3$-manifolds. The source reports substantial analytic groundwork but no complete general counting theory, and leaves several defining conventions unspecified.
\subsection{Sources}
\begin{itemize}
\item Local UnsolvedMath registry, permanent ID 2691, source code catalogue entry, ``Kirby Problem 1.32''; complete supplied three-part statement and literature notes.
\item \emph{K3: A New Problem List in Low-Dimensional Topology}, Problem 1.32 (2026), supplied source: \url{https://aimath.org/pastworkshops/kirbylistrep.pdf}. Source Problem 3.68 is a background cross-reference, not a permanent-ID identification.
\item [Kho00] Mikhail Khovanov, ``A categorification of the Jones polynomial,'' \emph{Duke Mathematical Journal} 101(3) (2000), 359--426, doi:10.1215/S0012-7094-00-10131-7.
\item [Wit12] Edward Witten, ``Fivebranes and knots,'' \emph{Quantum Topology} 3(1) (2012), 1--137, doi:10.4171/QT/26.
\item [KW07] Anton Kapustin and Edward Witten, ``Electric-magnetic duality and the geometric Langlands program,'' \emph{Communications in Number Theory and Physics} 1(1) (2007), 1--236, doi:10.4310/CNTP.2007.v1.n1.a1.
\item [GW12] Davide Gaiotto and Edward Witten, ``Knot invariants from four-dimensional gauge theory,'' \emph{Advances in Theoretical and Mathematical Physics} 16(3) (2012), 935--1086, doi:10.4310/atmp.2012.v16.n3.a5.
\item [Hay15] Andriy Haydys, ``Fukaya-Seidel category and gauge theory,'' \emph{Journal of Symplectic Geometry} 13(1) (2015), 151--207, doi:10.4310/JSG.2015.v13.n1.a5.
\item [Prz91] Józef H. Przytycki, ``Skein modules of 3-manifolds,'' \emph{Bulletin of the Polish Academy of Sciences, Mathematics} 39(1--2) (1991), 91--100.
\item [Tur88] V. G. Turaev, ``The Conway and Kauffman modules of a solid torus,'' \emph{Zap. Nauchn. Sem. Leningrad. Otdel. Mat. Inst. Steklov. (LOMI)} 167 (Issled. Topol. 6) (1988), 79--89, 190, doi:10.1007/BF01099241.
\item [GPV17] Sergei Gukov, Pavel Putrov, and Cumrun Vafa, ``Fivebranes and 3-manifold homology,'' \emph{Journal of High Energy Physics} 2017(7) (2017), 071, front matter+80, doi:10.1007/JHEP07(2017)071.
\item [GPPV20] Sergei Gukov, Du Pei, Pavel Putrov, and Cumrun Vafa, ``BPS spectra and 3-manifold invariants,'' \emph{Journal of Knot Theory and Its Ramifications} 29(2) (2020), 2040003, 85 pages, doi:10.1142/S0218216520400039.
\item [APS04] Marta M. Asaeda, Józef H. Przytycki, and Adam S. Sikora, ``Categorification of the Kauffman bracket skein module of I-bundles over surfaces,'' \emph{Algebraic and Geometric Topology} 4 (2004), 1177--1210, doi:10.2140/agt.2004.4.1177.
\item [Man07] Vassily Olegovich Manturov, ``Khovanov homology for virtual knots with arbitrary coefficients,'' \emph{Journal of Knot Theory and Its Ramifications} 16(3) (2007), 345--377, doi:10.1142/S0218216507005336.
\item [Gab13] Boštjan Gabrovšek, ``The categorification of the Kauffman bracket Skein module of $\mathbb{RP}^3$,'' \emph{Bulletin of the Australian Mathematical Society} 88(3) (2013), 407--422, doi:10.1017/S0004972713000105.
\item [Roz10] Lev Rozansky, ``A categorification of the stable SU(2) Witten-Reshetikhin-Turaev invariant of links in $S^2\times S^1$'' (2010), arXiv:1011.1958.
\item [Wil21] Michael Willis, ``Khovanov homology for links in'' connected sums of $S^2\times S^1$, \emph{Michigan Mathematical Journal} 70(4) (2021), 675--748, doi:10.1307/mmj/1594281620. The supplied title's connected-sum notation is garbled; its subject is identified in the accompanying background.
\item [MWW22] Scott Morrison, Kevin Walker, and Paul Wedrich, ``Invariants of 4-manifolds from Khovanov-Rozansky link homology,'' \emph{Geometry and Topology} 26(8) (2022), 3367--3420, doi:10.2140/gt.2022.26.3367.
\item [MW14] Rafe Mazzeo and Edward Witten, ``The Nahm pole boundary condition,'' in \emph{The influence of Solomon Lefschetz in geometry and topology}, Contemporary Mathematics 621 (2014), 171--226, American Mathematical Society, Providence, Rhode Island, doi:10.1090/conm/621/12422.
\item [MW20] Rafe Mazzeo and Edward Witten, ``The KW equations and the Nahm pole boundary condition with knots,'' \emph{Communications in Analysis and Geometry} 28(4) (2020), 871--942, doi:10.4310/CAG.2020.v28.n4.a4.
\item [Tau13] Clifford Henry Taubes, ``Compactness theorems for SL(2;C) generalizations of the 4-dimensional anti-self dual equations'' (2013), arXiv:1307.6447.
\item [LT20] Naichung Conan Leung and Ryosuke Takahashi, ``Energy bound for Kapustin-Witten solutions on $S^3\times\R_+$,'' \emph{International Mathematics Research Notices} 2020(19) (2020), 6135--6148, doi:10.1093/imrn/rny198.
\item [He19] Siqi He, ``A gluing theorem for the Kapustin-Witten equations with a Nahm pole,'' \emph{Journal of Topology} 12(3) (2019), 855--915, doi:10.1112/topo.12102.
\item [HM19a] Siqi He and Rafe Mazzeo, ``The extended Bogomolny equations and generalized Nahm pole boundary condition,'' \emph{Geometry and Topology} 23(5) (2019), 2475--2517, doi:10.2140/gt.2019.23.2475.
\item [Tau18] Clifford Henry Taubes, ``Sequences of Nahm pole solutions to the SU(2) KapustinWitten equations'' (2018), arXiv:1805.02773. The source singles out Theorem B's positive-Ricci-curvature restriction.
\item [Tau19] Clifford Henry Taubes, ``The R invariant solutions to the Kapustin Witten equations on $(0,8)\times\R^2\times\R$ with generalized Nahm pole asymptotics'' (2019), arXiv:1903.03539. The interval $(0,8)$ is retained as printed in the supplied title and may be extraction damage; it is not silently changed to an infinite interval.
\item Supplied literature triage checked 2026-08-17, OPEN-TRIAGE; it reports no completed analytic construction with all requested properties, records no source alteration, and cautions that dated triage does not exclude later or unindexed results.
\end{itemize}
\end{document}
